\documentclass{article}
\usepackage[utf8]{inputenc}
\usepackage{amsmath,amssymb}
\usepackage{graphicx}
\usepackage{hyperref}
\usepackage[margin=1in]{geometry}

\title{SmallGBM: Gradient Boosting with Bayesian Leaf Regularization\\ for Small-Sample Tabular Data}

\author{Ilya Emelyanov\\
\href{https://github.com/nsdmlk/SmallGBM}{github.com/nsdmlk/SmallGBM}\\
\texttt{nsdmlk@yandex.ru}}

\date{August 2026}

\begin{document}
\maketitle

\begin{abstract}
Gradient boosting machines (GBMs) are state-of-the-art for tabular data, yet leading implementations---XGBoost, LightGBM---are optimized for large-sample regimes and overfit severely on datasets with fewer than 1000 rows. We introduce SmallGBM, a gradient boosting library designed from first principles for small-sample tabular data. SmallGBM replaces standard leaf-weight estimation with a Bayesian shrinkage estimator that prevents overfitting on sparse leaves, and removes variance-inducing mechanisms like bootstrap sampling and column subsampling. We benchmark SmallGBM against XGBoost, LightGBM, and Random Forest on synthetic and real-world datasets with 20--500 training samples. SmallGBM shows superior stability under label noise (maintaining higher AUC at 10--25\% corruption) and near-identical performance to baselines on clean data. The library is pip-installable, scikit-learn compatible, and released under the MIT license.
\end{abstract}

\section{Introduction}
Gradient boosting machines \cite{friedman2001greedy} have become the default choice for tabular data, with XGBoost \cite{chen2016xgboost} and LightGBM \cite{ke2017lightgbm} dominating benchmarks. However, these libraries are engineered for datasets with thousands to millions of rows. Their default hyperparameters---deep trees, aggressive learning rates, subsampling---assume abundant data. On small datasets, these defaults lead to overfitting and high variance.

Small-sample problems are common in applied ML: medical trials with 50--200 patients, early-stage startup analytics with hundreds of users, and scientific experiments with limited observations. Practitioners working with such data often resort to linear models, losing the predictive power of GBMs.

SmallGBM fills this gap. It provides a gradient boosting implementation specifically designed for datasets with $n < 1000$, trading some peak performance on large data for stability and reliability when data is scarce.

\section{Method}
\label{sec:method}

SmallGBM is a gradient boosting classifier and regressor built on top of custom Bayesian decision trees. The algorithm follows the standard gradient boosting framework \cite{friedman2001greedy} with modifications specific to small-sample regimes.

\subsection{Bayesian Leaf Weight Estimation}
The core innovation is the leaf weight calculation. Standard gradient boosting computes leaf weights as the mean of pseudo-residuals:

\begin{equation}
w_{\text{leaf}} = \frac{\sum_{i \in L} r_i}{|L|}
\end{equation}

where $L$ is the set of samples in the leaf and $r_i$ are the pseudo-residuals. When $|L|$ is small (common with limited data and shallow trees), this estimate has high variance---a single outlier can dominate the weight.

SmallGBM uses a Bayesian shrinkage estimator:

\begin{equation}
w_{\text{leaf}} = \frac{\sum_{i \in L} r_i}{|L| + \sigma^2_{\text{noise}} / \sigma^2_{\text{prior}}}
\end{equation}

where $\sigma^2_{\text{noise}}$ is the variance of residuals in the leaf and $\sigma^2_{\text{prior}}$ is a hyperparameter controlling the strength of the prior belief that leaf weights should be close to zero. When $|L|$ is small, the denominator is dominated by the shrinkage term, pulling weights toward zero. As $|L|$ grows, the estimate converges to the standard mean.

\subsection{Design Choices for Small Data}
Several standard GBM techniques are intentionally disabled:

\begin{itemize}
    \item \textbf{No bootstrap}: Each tree is trained on the full dataset. Bootstrap sampling reduces effective sample size and adds variance when $n$ is already small.
    \item \textbf{No column subsampling}: All features are considered at each split. Feature subsampling is useful for high-dimensional data but harmful when both $n$ and $d$ are limited.
    \item \textbf{Shallow trees}: Default maximum depth of 3 with a minimum of 3 samples per leaf. Deep trees memorize noise on small datasets.
    \item \textbf{Adaptive prior}: The hyperparameter $\sigma_{\text{prior}}$ can automatically scale as $1/\sqrt{n}$.
\end{itemize}

\subsection{Implementation}
SmallGBM is implemented in Python with NumPy and scikit-learn. The library provides two main classes: \texttt{SmallGBMClassifier} and \texttt{SmallGBMRegressor}, both following the scikit-learn API (\texttt{fit}, \texttt{predict}, \texttt{predict\_proba}). The package is available via \texttt{pip install smallgbm}.

\section{Experiments}
\label{sec:experiments}

We evaluate SmallGBM on synthetic datasets with 20--150 training samples, comparing against XGBoost (v2.1), LightGBM (v4.5), and Random Forest (scikit-learn 1.5). All models use comparable hyperparameters (50 estimators, max\_depth=3, min\_samples\_leaf=3). Each experiment is repeated 10 times with different random splits. Error bars show $\pm 1$ standard deviation.

\subsection{Stability Under Label Noise}
We generate a 200-sample binary classification dataset and randomly flip 0--30\% of training labels. Figure~\ref{fig:noise} shows that SmallGBM degrades more gracefully than all baselines. At 20\% label noise, SmallGBM achieves the highest AUC, demonstrating the benefit of Bayesian regularization when labels are unreliable.

\begin{figure}[h]
\centering
\includegraphics[width=0.9\textwidth]{noise_comparison.png}
\caption{Test AUC under increasing label noise. SmallGBM (green) degrades slower than XGBoost, LightGBM, and Random Forest.}
\label{fig:noise}
\end{figure}

\subsection{Prior Sensitivity}
We vary $\sigma_{\text{prior}}$ from 0.01 to 100 on a 40-sample dataset. Figure~\ref{fig:prior} shows that performance is nearly flat across three orders of magnitude, indicating that the Bayesian prior is not a sensitive hyperparameter.

\begin{figure}[h]
\centering
\includegraphics[width=0.9\textwidth]{prior_sensitivity.png}
\caption{AUC vs. prior strength $\sigma_{\text{prior}}$. Performance is robust across two orders of magnitude.}
\label{fig:prior}
\end{figure}

\subsection{Learning Curve}
We train SmallGBM on increasing sample sizes from 15 to 150. Figure~\ref{fig:learning} shows a clear elbow at $n \approx 40$, beyond which performance stabilizes and variance decreases. Before this point, SmallGBM provides conservative predictions rather than overfitting.

\begin{figure}[h]
\centering
\includegraphics[width=0.9\textwidth]{learning_curve.png}
\caption{Test AUC as a function of training set size. Reliable performance begins at $n \approx 40$.}
\label{fig:learning}
\end{figure}

\subsection{Regression Performance}
SmallGBMRegressor shows the same stability properties as the classifier. RMSE decreases with sample size, with the elbow at similar $n \approx 40-60$ range.

\subsection{Class Imbalance}
On datasets with minority class proportions from 5\% to 50\%, SmallGBM maintains AUC above 0.7 even at extreme imbalance. However, XGBoost and LightGBM are more robust to severe imbalance, suggesting that Bayesian shrinkage can be too conservative for rare class detection.

\section{Discussion}
SmallGBM is not a universal improvement over existing GBM libraries. On clean, balanced datasets with adequate sample size, XGBoost and LightGBM match or slightly exceed its accuracy. However, SmallGBM offers distinct advantages:

\begin{itemize}
    \item \textbf{Noise resilience}: Superior performance when labels are corrupted.
    \item \textbf{Stability}: Lower variance across random seeds and data splits.
    \item \textbf{Predictability}: Smooth learning curve with no catastrophic failures.
    \item \textbf{Simplicity}: Sensible defaults for small data with no hyperparameter tuning required.
\end{itemize}

These properties make SmallGBM a safe choice for practitioners working with limited, potentially noisy data who cannot afford extensive hyperparameter tuning.

\subsection{Limitations}
The current version has several limitations: (1) native support only for numeric features, (2) class imbalance handling trails dedicated libraries, (3) no GPU acceleration, (4) training time scales with the full dataset (no subsampling).

\section{Conclusion and Future Work}
SmallGBM demonstrates that gradient boosting can be redesigned for small-sample regimes by incorporating Bayesian regularization and removing variance-inducing mechanisms. The library is production-ready and available on PyPI.

Future directions include: automatic categorical encoding (target encoding with internal cross-validation), adaptive prior selection based on meta-features of the dataset, built-in feature selection for high-dimensional small-$n$ problems, and integration with the MetaLearn framework for automatic model selection based on dataset characteristics.

\bibliographystyle{unsrt}
\bibliography{refs}

\end{document}